% Lösungen Einheit 3 von 4 – Lage als Anhang (zusammenbau v0.1)
\einheitenkopf{Einheit 3 von 4 · Lage als Anhang}

\erg{16}{a) ob ein Stützpunkt von $g$ auf $h$ liegt – er entscheidet zwischen identisch und echt parallel \quad b) $(3 | 1 | 0)$ ist kein Vielfaches von $(1 | 3 | 0)$: die erste Koordinate verlangt den Faktor 3, die zweite den Faktor $\tfrac{1}{3}$ – verschiedene Richtungen; $g$ und $h$ schneiden sich im gemeinsamen Stützpunkt und sind nicht identisch. \quad c) $\overrightarrow{PQ} = (4 | -2 | 5)$ ist kein Vielfaches von $(2 | -1 | 3)$ (Faktor 2 in den ersten beiden Koordinaten, $\tfrac{5}{3}$ in der dritten) – $Q$ liegt nicht auf $g$, die Geraden sind echt parallel und nicht identisch. \quad d) $\overrightarrow{AB} = (30 | 20 | 3)$, $\overrightarrow{CD} = (60 | 40 | 6) = 2 \cdot \overrightarrow{AB}$ – die Tunnelachsen sind parallel. \quad e) $h\colon \vec x = (5 | 5 | 0) + t \cdot (-5 | 0 | 5)$; $AC$ liegt in der $x_1x_2$-Ebene, $h$ durchstößt diese Ebene nur in $B$, und $B$ liegt nicht auf $AC$; die Richtungen $(-5 | 5 | 0)$ und $(-5 | 0 | 5)$ sind nicht parallel – windschief. \quad f) Beide Geraden liegen in $E$; die Richtungen $(2 | 2 | 4)$ und $(2 | 2 | -2)$ sind keine Vielfachen – zwei nicht parallele Geraden einer Ebene schneiden sich. \quad g) Nein. $g$ und $p$ liegen in einer Ebene, eine senkrecht schneidende Gerade muss nicht in ihr liegen: $q\colon \vec x = (1 | 3 | 1) + t \cdot (1 | 0 | -2)$ schneidet $g$ senkrecht ($2 + 0 - 2 = 0$), hat aber stets die zweite Koordinate 3, $p$ stets 5 – $q$ und $p$ sind windschief.}
\erg{17}{a) Ein gemeinsamer Punkt genügt nicht: $(2 | 1 | 1)$ ist kein Vielfaches von $(1 | 2 | -1)$, die Geraden schneiden sich in $(4 | -1 | 2)$ und sind nicht identisch. \quad b) In einer Ebene sind zwei Geraden entweder parallel oder schneidend; windschief ist nur möglich, wenn es keine gemeinsame Ebene gibt.}
